\chapter{Sistemas Metabólicos}
\label{cap:08}

O metabolismo é o conjunto de reações bioquímicas que mantém a vida: converte nutrientes
em energia (ATP), em blocos construtores (aminoácidos, nucleotídeos, lipídios) e em
biomassa. Em uma célula de \textit{Escherichia coli} crescendo em glicose, mais de 2.500
reações catalisadas por enzimas operam simultaneamente, interconectadas numa rede cuja
topologia determina quais fenótipos metabólicos são possíveis. Compreender essa rede como
sistema --- não apenas reação por reação, mas em sua integridade estrutural e dinâmica ---
é o tema deste capítulo.

Abordaremos o problema em quatro escalas progressivamente mais abstratas: cinética
individual de reações enzimáticas; análise de fluxo metabólico (MFA) para determinar
experimentalmente os fluxos numa rede; análise de balanço de fluxo (FBA) para predizer
computacionalmente distribuições de fluxo em redes de escala genômica; e modos elementares
de fluxo como representação estrutural do espaço de soluções possíveis. Encerramos com
a metabolômica e sua integração com os demais níveis ômicos.

\begin{figure}[htbp]
\centering
\begin{tikzpicture}[node distance=2.5cm, scale=0.95, transform shape]
    \node[draw, circle, fill=laranjacel!30, text width=2cm, align=center] (nut) {Nutrientes};
    \node[draw, circle, fill=azulclaro!30, text width=2cm, align=center, right=of nut] (met) {Metabolismo};
    \node[draw, circle, fill=verdeacido!30, text width=1.5cm, align=center, above right=1cm and 1.5cm of met] (atp) {Energia\\(ATP)};
    \node[draw, circle, fill=roxodna!30, text width=1.5cm, align=center, below right=1cm and 1.5cm of met] (bio) {Biomassa};

    \draw[->, thick, azulescuro, line width=1.5pt] (nut) -- (met);
    \draw[->, thick, azulescuro, line width=1.5pt] (met) -- (atp);
    \draw[->, thick, azulescuro, line width=1.5pt] (met) -- (bio);
\end{tikzpicture}
\caption{Visão esquemática do metabolismo celular. Nutrientes externos (glicose, aminoácidos, etc.) entram na rede metabólica, onde são processados por centenas a milhares de reações enzimáticas. Os produtos finais bifurcados são a \emph{energia química} (ATP, NADH --- para funções celulares) e a \emph{biomassa} (proteínas, lipídios, ácidos nucleicos --- para crescimento e divisão).}
\label{fig:08-metabolismo}
\end{figure}

\section{Cinéticas de Reações Metabólicas}

\subsection{Lei de Ação de Massas e Michaelis-Menten Reversível}

A cinética química mais fundamental é a \emph{lei de ação de massas}: para uma reação
elementar $A + B \xrightarrow{k} C$, a velocidade é proporcional ao produto das
concentrações dos reagentes: $v = k[A][B]$. Para a reação unimolecular $A \xrightarrow{k}
B$, $v = k[A]$, com solução $[A](t) = [A]_0 e^{-kt}$. Essa lei aplica-se estritamente a
reações elementares --- sem intermediários catalíticos.

A cinética de Michaelis-Menten (discutida no Cap.\ 7) é o modelo padrão para catálise
enzimática com um substrato. Para redes metabólicas com reações reversíveis --- a maioria,
pois poucas reações são termodinamicamente irreversíveis sob condições celulares --- é
necessária a forma reversível:

\begin{equation}
v = \frac{V_f \frac{[S]}{K_{MS}} - V_r \frac{[P]}{K_{MP}}}{1 + \frac{[S]}{K_{MS}} + \frac{[P]}{K_{MP}}}
\end{equation}
A consistência termodinâmica exige $K_{eq} = V_f K_{MP} / (V_r K_{MS})$ --- a constante
de equilíbrio é determinada pelos parâmetros cinéticos, não independentemente.

\subsection{Convenience Kinetics e S-systems}

As \emph{convenience kinetics} (Liebermeister e Klipp, 2006) generalizam o formalismo MM
para reações com múltiplos substratos e produtos, incorporando inibidores e ativadores
de forma termodinamicamente consistente. São amplamente usadas em modelos cinéticos
detalhados de vias metabólicas centrais.

Os \emph{S-systems} (Savageau, 1969) adotam uma abordagem diferente: aproximam cada taxa
metabólica por uma lei de potência $v = \alpha \prod_j X_j^{g_j}$, separando síntese de
degradação:

\begin{equation}
\frac{dX_i}{dt} = \alpha_i \prod_{j=1}^{n} X_j^{g_{ij}} - \beta_i \prod_{j=1}^{n} X_j^{h_{ij}}
\end{equation}
Os expoentes $g_{ij}$ e $h_{ij}$ são as elasticidades cinéticas, equivalentes aos
coeficientes de sensibilidade da MCA. A grande vantagem é analítica: o estado estacionário
é encontrado por um sistema de equações lineares em escala logarítmica, permitindo análise
algébrica de estabilidade e sensibilidade. Essa é a base da \emph{Biochemical Systems
Theory} (BST) de Voit e Savageau.

\begin{lstlisting}[language=Python, caption=Via metabolica com cinetica Michaelis-Menten]
import numpy as np
from scipy.integrate import odeint
import matplotlib.pyplot as plt

def via_metabolica(y, t, params):
    # Via: A -> B -> C, com ramificacao B -> D
    A, B, C, D = y
    Vmax1, Km1, Vmax2, Km2, Vmax3, Km3 = params
    v1 = Vmax1 * A / (Km1 + A)   # A -> B
    v2 = Vmax2 * B / (Km2 + B)   # B -> C
    v3 = Vmax3 * B / (Km3 + B)   # B -> D
    return [-v1, v1 - v2 - v3, v2, v3]

params = [10.0, 1.0, 6.0, 1.5, 4.0, 2.0]
y0 = [50.0, 0.0, 0.0, 0.0]
t = np.linspace(0, 30, 500)
sol = odeint(via_metabolica, y0, t, args=(params,))

for i, label in enumerate(['A', 'B', 'C', 'D']):
    plt.plot(t, sol[:, i], label=label, linewidth=2)
plt.xlabel('Tempo (min)'); plt.ylabel('[Metabolito] (mM)')
plt.legend(); plt.grid(True); plt.tight_layout()
\end{lstlisting}


\section{Análise de Fluxo Metabólico (MFA)}

\begin{definicaoenv}[Análise de Fluxo Metabólico]
A MFA (\textit{Metabolic Flux Analysis}) é uma técnica experimental e computacional para
determinar as taxas de reação (fluxos metabólicos) numa rede, combinando balanços de massa
em estado estacionário com medições experimentais de taxas de consumo e produção de
metabólitos.
\end{definicaoenv}

\subsection{Representação Matricial: a Matriz Estequiométrica}

A topologia de uma rede metabólica é codificada na \emph{matriz estequiométrica} $S$,
de dimensão $m \times n$, onde $m$ é o número de metabólitos e $n$ o número de reações.
O elemento $S_{ij}$ é o coeficiente estequiométrico do metabólito $i$ na reação $j$:
positivo se produzido, negativo se consumido, zero se ausente.

Para a rede simples $v_1: A \to B$, $v_2: B \to C$, $v_3: B \to D$:

\begin{equation}
S = \begin{bmatrix}
-1 & 0 & 0 \\
1 & -1 & -1 \\
0 & 1 & 0 \\
0 & 0 & 1
\end{bmatrix}
\begin{matrix} \leftarrow A \\ \leftarrow B \\ \leftarrow C \\ \leftarrow D \end{matrix}
\end{equation}
O balanço dinâmico de massa para todos os metabólitos é simplesmente:
\begin{equation}
\frac{d\mathbf{c}}{dt} = S \cdot \mathbf{v}
\end{equation}
onde $\mathbf{c}$ é o vetor de concentrações e $\mathbf{v}$ o vetor de fluxos. A hipótese
fundamental da MFA é a do \emph{pseudo-estado estacionário}: os metabólitos intracelulares
atingem estado estacionário muito mais rapidamente do que as variáveis externas (biomassa,
nutrientes no meio), portanto $d\mathbf{c}/dt \approx \mathbf{0}$ e:

\begin{equation}
S \cdot \mathbf{v} = \mathbf{0}
\end{equation}
Esse sistema de equações lineares homogêneas é geralmente \emph{subdeterminado}: há mais
reações ($n$) do que metabólitos ($m$), portanto o espaço nulo de $S$ é não-trivial e
infinitas distribuições de fluxo satisfazem o balanço de massa. Para determinar $\mathbf{v}$
uniquamente, são necessárias medições adicionais de fluxos de troca (taxas de consumo de
glicose, de produção de CO$_2$, de secreção de metabólitos) ou dados de marcação isotópica.

\begin{lstlisting}[language=Python, caption=Resolucao de fluxos por balanco de massa]
import numpy as np
from scipy.linalg import lstsq

# Matriz estequiometrica (metabolitos A,B,C,D x reacoes v1..v6)
S = np.array([
    [ 1, -1,  0,  0,  0,  0],  # A
    [ 0,  1, -1, -1,  0,  0],  # B
    [ 0,  0,  1,  0, -1,  0],  # C
    [ 0,  0,  0,  1,  0, -1]   # D
])

# Fluxos medidos: v1=10 (uptake A), v5=6 (producao C), v6=4 (producao D)
meas_idx   = [0, 4, 5]   # indices das reacoes medidas
unmeas_idx = [1, 2, 3]   # indices das reacoes nao-medidas

v_meas = np.array([10.0, 6.0, 4.0])
b = -S[:, meas_idx] @ v_meas

v_unmeas, _, _, _ = lstsq(S[:, unmeas_idx], b)

v_all = np.zeros(6)
v_all[meas_idx]   = v_meas
v_all[unmeas_idx] = v_unmeas

print("Distribuicao de fluxos:")
for i, flux in enumerate(v_all, 1):
    print(f"  v{i} = {flux:.2f}")

# Verificar balanco: S*v deve ser ~0
residual = S @ v_all
print(f"Residual max: {np.abs(residual).max():.2e}")
\end{lstlisting}

\subsection{$^{13}$C-MFA: Rastreamento Isotópico}

A MFA baseada em marcação com $^{13}$C (\textit{isotope labeling MFA}) resolve ambiguidades
que surgem quando a rede tem mais graus de liberdade do que medições de troca. Células são
cultivadas com substrato marcado --- por exemplo, [1,2-$^{13}$C]-glicose --- e os padrões
de distribuição de marcação nos metabólitos intracelulares são medidos por espectrometria
de massa ou NMR. Como a distribuição de $^{13}$C nos metabólitos depende diretamente dos
fluxos pelos quais os carbonos passaram, ela contém informação sobre fluxos em vias
paralelas, a direção de reações reversíveis e a extensão de ciclos metabólicos --- informações
inacessíveis apenas pelos balanços de massa externos.

O procedimento envolve: (1) definir a topologia da rede e mapear as transições atômicas
de carbono em cada reação; (2) simular, para um dado conjunto de fluxos, o padrão de
marcação esperado nos metabólitos; (3) ajustar os fluxos para minimizar a diferença
(SSE ponderado) entre os padrões simulado e experimental, usando algoritmos de estimação
de parâmetros como os discutidos no Capítulo 5. A ferramenta padrão para $^{13}$C-MFA
é o Iso2Flux ou o OpenFlux.


\section{Análise de Balanço de Fluxo (FBA)}

\begin{definicaoenv}[Flux Balance Analysis]
FBA (\textit{Flux Balance Analysis}) é um método computacional baseado em programação linear
para predizer distribuições de fluxo metabólico em redes de escala genômica, sem requerer
parâmetros cinéticos. Assume estado estacionário e maximiza (ou minimiza) uma função
objetivo biológica sujeita às restrições estequiométricas e de capacidade.
\end{definicaoenv}

\subsection{Formulação Matemática}

O problema de FBA é formulado como programação linear:

\begin{align}
\text{maximizar} \quad & \mathbf{c}^T \mathbf{v} \notag \\
\text{sujeito a} \quad & S \cdot \mathbf{v} = \mathbf{0} \\
& \mathbf{v}_{min} \leq \mathbf{v} \leq \mathbf{v}_{max} \notag
\end{align}
O vetor $\mathbf{c}$ define o objetivo: tipicamente $c_i = 1$ para a reação de biomassa e
$c_i = 0$ para todas as demais, equivalendo a maximizar a taxa de crescimento. Os limites
$\mathbf{v}_{max}$ e $\mathbf{v}_{min}$ codificam restrições termodinâmicas (reações
irreversíveis têm $v_{min} = 0$), de capacidade enzimática e de disponibilidade de
nutrientes (limitações de uptake).

A justificativa biológica para a função objetivo de biomassa é evolutiva: microrganismos
submetidos a seleção para crescimento máximo no nicho ecológico relevante devem ter
evoluído para distribuições de fluxo próximas ao ótimo do FBA. Essa hipótese, validada
extensivamente em \textit{E. coli}, \textit{S. cerevisiae} e outros organismos, explica
por que o FBA prevê taxas de crescimento e distribuições de fluxo experimentalmente
próximas do observado, sem qualquer parâmetro cinético.

\subsection{Modelos de Escala Genômica (GEMs)}

Os \emph{genome-scale metabolic models} (GEMs) são reconstruções computacionais de todo
o metabolismo de um organismo, baseadas na anotação do genoma e na literatura bioquímica.
O GEM de \textit{E. coli} iJO1366 contém 2.583 reações e 1.805 genes; o modelo humano
Recon3D tem 13.543 reações e 3.288 genes. Esses modelos estão disponíveis em repositórios
públicos como BiGG Models (bigg.ucsd.edu) e no formato SBML para importação em ferramentas
como COBRApy.

\subsection{FVA e Predição de Essencialidade Gênica}

A \emph{análise de variabilidade de fluxo} (FVA) determina, para cada reação, o intervalo
$[v_{\min}^{opt}, v_{\max}^{opt}]$ de valores que essa reação pode assumir enquanto a
função objetivo permanece próxima ao ótimo (tipicamente $\geq 90\%$ do valor máximo). FVA
identifica reações essencialmente fixas (intervalo estreito), que provavelmente carregam
fluxo obrigatório, e reações com alta flexibilidade --- indicando redundância metabólica.

Para predição de essencialidade gênica, o FBA é resolvido repetidamente com cada gene
deletado in silico: a reação (ou conjunto de reações) codificada pelo gene é bloqueada
($v_{max} = 0$) e o crescimento resultante é comparado ao tipo selvagem. Um gene é
classificado como essencial se sua deleção reduz o crescimento a zero; essencial
condicionalmente se reduz apenas em certos meios; e dispensável se o crescimento não
é afetado (indicando rotas alternativas). O FBA prevê corretamente $\sim$85--90\% dos
genes essenciais experimentalmente determinados em \textit{E. coli} em meio de glicose
aeróbico.

\begin{lstlisting}[language=Python, caption=FBA basico com COBRApy]
import cobra
from cobra.io import load_model
from cobra.flux_analysis import flux_variability_analysis, single_gene_deletion

# Carregar GEM de E. coli
model = load_model("iJO1366")
print(f"Reacoes: {len(model.reactions)} | Genes: {len(model.genes)}")

# Definir meio de cultura (glicose aerobica)
model.medium = {"EX_glc__D_e": 10.0, "EX_o2_e": 20.0,
                "EX_nh4_e": 1000.0, "EX_pi_e": 1000.0}

# FBA
sol = model.optimize()
print(f"Taxa de crescimento: {sol.objective_value:.3f} 1/h")
print(f"Uptake de O2: {sol.fluxes['EX_o2_e']:.2f}")
print(f"Producao de CO2: {sol.fluxes['EX_co2_e']:.2f}")

# Flux Variability Analysis (reacoes da glicolise central)
glycolysis = ["PGI", "PFK", "FBA", "TPI", "GAPD", "PGK", "PGM", "ENO", "PYK"]
fva = flux_variability_analysis(model, reaction_list=glycolysis,
                                 fraction_of_optimum=0.9)
print("\nFVA - Glicolise:")
for rxn, row in fva.iterrows():
    print(f"  {rxn}: [{row['minimum']:.2f}, {row['maximum']:.2f}]")

# Essencialidade genica (subset)
with model:
    del_results = single_gene_deletion(model)
essential = del_results[del_results["growth"] < 0.01]
print(f"\nGenes essenciais: {len(essential)}/{len(model.genes)}")
\end{lstlisting}


\section{Modos Elementares de Fluxo}

\begin{definicaoenv}[Modos Elementares de Fluxo]
Os \emph{modos elementares de fluxo} (EFMs, \textit{Elementary Flux Modes}) são os
conjuntos mínimos de reações que podem operar em estado estacionário, respeitando as
restrições de reversibilidade. São os ``caminhos funcionalmente independentes'' da rede:
cada EFM corresponde a uma rota metabólica completa de substrato a produto. Qualquer
distribuição de fluxo fisiologicamente realizável pode ser expressa como combinação linear
não-negativa de EFMs.
\end{definicaoenv}

Para a rede simples com dois caminhos paralelos de $A$ a $D$ ($A \to B \to D$ e
$A \to C \to D$), existem exatamente dois EFMs: aquele que usa o caminho superior
($v_2 = v_4 = 1$, $v_3 = v_5 = 0$) e aquele que usa o inferior ($v_3 = v_5 = 1$,
$v_2 = v_4 = 0$). Qualquer distribuição de fluxo fisiológica é uma mistura convexa
desses dois.

As \emph{vias extremas} (\textit{Extreme Pathways}) são um subconjunto dos EFMs que
formam a base convexa do cone de fluxos --- os vetores geradores a partir dos quais
qualquer fluxo viável pode ser construído. Matematicamente, são os geradores extremos
do poliedro definido por $S \cdot \mathbf{v} = 0$ e $\mathbf{v} \geq 0$ (para reações
irreversíveis).

O cálculo de EFMs enfrenta um desafio computacional severo: o número de EFMs cresce
exponencialmente com o tamanho da rede. Para o metabolismo central de \textit{E. coli}
($\sim$100 reações), existem dezenas de milhares de EFMs; para redes de escala genômica,
o número é astronomicamente grande e o cálculo torna-se inviável. Algoritmos como o
Double Description Method e o método de árvores de padrões de bits tornaram o cálculo
prático para redes de até $\sim$100 reações. Ferramentas como efmtool (ETH Zürich) e
CellNetAnalyzer implementam esses algoritmos.

As aplicações dos EFMs incluem: identificação de rotas metabólicas alternativas para
um objetivo de engenharia; análise de robustez (quantas rotas são eliminadas por knockout
de um gene); análise de optimalidade (comparar o fluxo observado com a mistura de EFMs
mais eficiente); e identificação de alvos de knockout sintético letal (pares de reações
cuja remoção simultânea elimina todos os EFMs viáveis).

\section{Metabolômica}

\begin{definicaoenv}[Metabolômica]
Metabolômica é o estudo abrangente do conjunto completo de metabólitos (metaboloma) numa
amostra biológica. Enquanto o genoma e o transcriptoma descrevem o potencial e a
intenção celular, o metaboloma reflete o fenótipo bioquímico atual --- o resultado
integrado de todas as regulações genéticas, transcricionais, proteicas e ambientais.
\end{definicaoenv}

O metaboloma humano compreende $\sim$114.000 compostos catalogados no Human Metabolome
Database (HMDB, hmdb.ca), dos quais $\sim$2.500 são detectáveis endogenamente por
técnicas analíticas modernas. Os metabólitos abrangem uma enorme faixa de propriedades
químicas --- de metabólitos polares como a glicose (180 Da) a lipídios apolares como
o colesterol (387 Da) e a ceramidas ($>$700 Da) --- o que exige múltiplas plataformas
analíticas para cobertura abrangente.

\subsection{Plataformas Analíticas}

A \emph{cromatografia líquida acoplada a espectrometria de massa} (LC-MS) é a plataforma
mais versátil, combinando separação cromatográfica com detecção por massa. Em modo de
ionização por electrospray (ESI), funciona para metabólitos polares e apolares. Em modo
tandem (MS/MS), fragmenta os metabólitos selecionados para confirmação estrutural e
identificação de novos compostos.

A \emph{cromatografia gasosa} (GC-MS) oferece altíssima reprodutibilidade e extensas
bibliotecas espectrais (NIST, Wiley), mas requer derivatização química prévia para
tornar os metabólitos voláteis, limitando sua aplicação a compostos que toleram esse
processamento (aminoácidos, ácidos orgânicos, açúcares simples).

A \emph{ressonância magnética nuclear} (NMR) é não-destrutiva, quantitativa e não requer
separação cromatográfica, mas tem sensibilidade $\sim$1000 vezes menor que a LC-MS.
É particularmente útil para identificação estrutural de metabólitos desconhecidos (por
$^{1}$H e $^{13}$C NMR bidimensional) e para $^{13}$C-MFA.

\subsection{Targeted vs. Untargeted}

A \emph{metabolômica targeted} foca num conjunto predefinido de metabólitos, geralmente
quantificados por padrões de referência e monitoramento de transições específicas de MS
(SRM/MRM). Produz dados quantitativos absolutos e altamente reprodutíveis, mas por
definição não pode descobrir novos biomarcadores.

A \emph{metabolômica untargeted} (ou global) busca detectar todos os metabólitos
ionizáveis na amostra sem seleção prévia. Produz milhares de \emph{features} (pares
massa/tempo de retenção), mas a maioria permanece não-identificada (nível MSI 4). O
principal desafio é a identificação: apenas $\sim$5--30\% dos features em estudos
untargeted são anotados com confiança. Bases de dados como HMDB, METLIN e mzCloud
fornecem espectros de referência para matching.

\subsection{Pipeline de Análise}

O pipeline padrão de metabolômica untargeted compreende: (1) preparação das amostras
com quenching metabólico rápido (para parar o metabolismo instantaneamente) e extração
líquida; (2) aquisição LC-MS com amostras QC intercaladas para monitorar deriva
instrumental; (3) detecção de picos e alinhamento de tempos de retenção (ferramentas:
XCMS, MZmine, MS-DIAL); (4) normalização (por padrão interno, por mediana ou por
proteína total); (5) análise estatística multivariada (PCA para exploração, PLS-DA para
classificação supervisionada) e univariada (teste $t$ com correção FDR); (6) análise
de enriquecimento de vias (MetaboAnalyst, mummichog).

\begin{lstlisting}[language=Python, caption=Analise exploratoria de dados metabolomicos]
import pandas as pd
import numpy as np
from sklearn.preprocessing import StandardScaler
from sklearn.decomposition import PCA
import matplotlib.pyplot as plt

# Simular dados: 20 amostras (10 ctrl + 10 trt), 200 metabolitos
np.random.seed(42)
n_ctrl, n_trt, n_met = 10, 10, 200
ctrl = np.random.lognormal(3, 1, (n_ctrl, n_met))
trt  = np.random.lognormal(3, 1, (n_trt,  n_met))
# 20 metabolitos alterados no tratamento
trt[:, :20] *= np.random.uniform(2, 5, 20)
X = np.vstack([ctrl, trt])
labels = ['ctrl']*n_ctrl + ['trt']*n_trt

# Pre-processamento: log2 + z-score
X_log = np.log2(X + 1)
scaler = StandardScaler()
X_scaled = scaler.fit_transform(X_log)

# PCA
pca = PCA(n_components=2)
scores = pca.fit_transform(X_scaled)
print(f"PC1: {pca.explained_variance_ratio_[0]*100:.1f}% | "
      f"PC2: {pca.explained_variance_ratio_[1]*100:.1f}%")

# Score plot
fig, ax = plt.subplots(figsize=(8, 6))
for group, color, marker in [('ctrl', 'blue', 'o'), ('trt', 'red', '^')]:
    idx = [i for i, l in enumerate(labels) if l == group]
    ax.scatter(scores[idx, 0], scores[idx, 1],
               c=color, marker=marker, label=group, s=80, alpha=0.8)
ax.set_xlabel(f'PC1 ({pca.explained_variance_ratio_[0]*100:.1f}%)')
ax.set_ylabel(f'PC2 ({pca.explained_variance_ratio_[1]*100:.1f}%)')
ax.set_title('PCA Score Plot - Metabolomica')
ax.legend(); ax.grid(True, alpha=0.3); plt.tight_layout()
\end{lstlisting}

\subsection{Integração Multi-Ômica}

A metabolômica ganha seu pleno poder quando integrada aos demais níveis ômicos. A
integração genômica-transcriptômica-proteômica-metabolômica (multi-omics) permite
conectar alterações genéticas (variantes, expressão diferencial) a consequências
bioquímicas mensuráveis. Ferramentas como o mapa de KEGG e o Reactome permitem mapear
metabólitos diferenciados em vias metabólicas, identificando quais vias estão ativadas
ou reprimidas num dado contexto.

A integração com GEMs é particularmente poderosa: dados de expressão de RNA (ou proteína)
podem ser usados para contextualizar o GEM --- restringindo reações cujos genes não são
expressos --- e produzir modelos específicos de tecido ou condição que refletem melhor
o metabolismo do sistema estudado. Algoritmos como GIMME, iMAT e FASTCORE implementam
essa integração.


\section{Exercícios}

\begin{exercicio}
Uma enzima catalisa $S \to P$ com $K_M = 5$ mM e $V_{\max} = 100$ $\mu$mol/min.
(a) Calcule a velocidade inicial quando $[S] = 10$ mM. (b) Qual $[S]$ é necessário para
atingir $v = 75\%$ de $V_{\max}$? (c) Se a concentração total de enzima dobrar, como
mudam $K_M$ e $V_{\max}$? (d) Simule numericamente a depleção de 50 mM de substrato ao
longo do tempo, incluindo degradação linear do substrato com $k_{deg} = 0.1$ min$^{-1}$.
\end{exercicio}

\begin{exercicio}[Balanço de Massa]
Para a rede com quatro metabólitos e seis reações: $v_1: \to A$, $v_2: A \to B$,
$v_3: B \to C$, $v_4: B \to D$, $v_5: C \to$, $v_6: D \to$: (a) escreva a matriz
estequiométrica $S$; (b) escreva as equações de balanço de massa em estado estacionário;
(c) dados $v_1 = 10$, $v_5 = 6$ e $v_6 = 4$, determine os fluxos desconhecidos por
resolução do sistema linear; (d) verifique que $S \cdot \mathbf{v} = \mathbf{0}$.
\end{exercicio}

\begin{exercicio}[FBA]
Considere um sistema metabólico simplificado: glicose ($v_{glc}$) pode ser metabolizada
por respiração ($v_{resp}$, rendimento 38 ATP/glc) ou fermentação ($v_{ferm}$, rendimento
2 ATP/glc). A biomassa requer 40 ATP/gDW. O uptake máximo de glicose é 10 mmol/gDW/h e
o sistema é aeróbico. (a) Formule o problema de FBA (variáveis, função objetivo, restrições).
(b) Qual é a taxa máxima de crescimento com respiração pura? (c) E com fermentação pura?
(d) Como o modelo se comportaria em anaerobiose (sem O$_2$)?
\end{exercicio}

\begin{exercicio}[S-systems]
Para o S-system unidimensional $\dot{X} = \alpha X_0^{g_0} - \beta X^h$, onde $X_0$ é o
substrato fixo: (a) encontre o estado estacionário $X^*$; (b) calcule a sensibilidade
logarítmica $\partial \ln X^* / \partial \ln \alpha$; (c) compare com o resultado da MCA
para o mesmo sistema: $S_{X^*, \alpha} = 1$ e $S_{X^*, \beta} = -1$ (indep. de $h$).
\end{exercicio}

\begin{exercicio}[COBRApy]
Usando COBRApy com o modelo iJO1366 de \textit{E. coli}: (a) simule o crescimento em
glicose aeróbica e anaeróbica; compare as taxas de crescimento e os fluxos de uptake;
(b) realize FVA com 90\% do ótimo para as reações da glicolise central e identifique
quais têm fluxo obrigatório (intervalo estreito) vs. flexível; (c) identifique os 10 genes
essenciais com maior impacto no crescimento (maior redução percentual quando deletados).
\end{exercicio}

\begin{exercicio}[Metabolômica]
Com um conjunto de dados metabolômicos simulado (controle vs. tratamento, 10 amostras
cada, 200 metabólitos): (a) aplique o pipeline completo de pré-processamento (log2 +
z-score + imputação de valores ausentes pela mediana); (b) construa um PCA score plot e
interprete a separação entre grupos; (c) identifique os 20 metabólitos mais significativos
(menor $p$-valor no teste $t$ com correção FDR); (d) esboce como você mapearia esses
metabólitos em vias do KEGG usando MetaboAnalyst.
\end{exercicio}

\begin{exercicio}[Projeto Integrador]
Escolha um organismo modelo (sugestão: \textit{Saccharomyces cerevisiae}) e: (a) baixe
o GEM Yeast8 do BiGG Models; (b) simule o crescimento em glicose aeróbica e anaeróbica;
(c) identifique os genes essenciais em cada condição e calcule a sobreposição; (d) realize
FVA para o metabolismo central do carbono e identifique reações com alternativas metabólicas;
(e) projete uma estratégia de engenharia metabólica para maximizar a produção de etanol
mantendo pelo menos 10\% do crescimento tipo selvagem. Interprete os resultados em termos
de biologia do organismo.
\end{exercicio}

\section{Notas Bibliográficas}

O livro de Palsson \cite{palsson2015} (\textit{Systems Biology: Constraint-based Reconstruction
and Analysis}) é a referência definitiva para FBA e GEMs. O artigo de Orth, Thiele e Palsson
\cite{orth2010} no \textit{Nature Biotechnology} é uma introdução acessível e amplamente
citada ao FBA. O COBRApy \cite{ebrahim2013} é a implementação Python de referência para
análise de GEMs.

Para MFA e $^{13}$C-MFA, Wiechert \cite{wiechert2001} é a referência metodológica clássica;
Antoniewicz (2015) \cite{antoniewicz2015} revisa os avanços mais recentes. Os S-systems
foram introduzidos por Savageau (1969) \cite{savageau1969} e amplamente desenvolvidos em Voit
\cite{voit2013}. Os modos elementares de fluxo foram formalizados por Schuster et al.\
\cite{schuster1999}; as vias extremas por Schilling et al.\ (2000).

Para metabolômica, Patti, Yanes e Siuzdak \cite{patti2012} fornecem uma revisão excelente.
A plataforma MetaboAnalyst \cite{chong2019} e o HMDB \cite{wishart2022} são os recursos
computacionais mais utilizados. Para integração multi-ômica com GEMs, Opdam et al.\ (2017)
e Lewis et al.\ (2012) \cite{lewis2012} são referências centrais.

\begin{thebibliography}{99}

\bibitem{palsson2015}
Palsson B.O. (2015).
\textit{Systems Biology: Constraint-based Reconstruction and Analysis}.
Cambridge University Press.

\bibitem{orth2010}
Orth J.D., Thiele I., Palsson B.O. (2010).
What is flux balance analysis?
\textit{Nature Biotechnology} \textbf{28}: 245--248.

\bibitem{ebrahim2013}
Ebrahim A. et al. (2013).
COBRApy: COnstraints-Based Reconstruction and Analysis for Python.
\textit{BMC Systems Biology} \textbf{7}: 74.

\bibitem{wiechert2001}
Wiechert W. (2001).
$^{13}$C metabolic flux analysis.
\textit{Metabolic Engineering} \textbf{3}: 195--206.

\bibitem{antoniewicz2015}
Antoniewicz M.R. (2015).
Methods and advances in metabolic flux analysis: a mini-review.
\textit{Journal of Industrial Microbiology and Biotechnology} \textbf{42}: 317--325.

\bibitem{savageau1969}
Savageau M.A. (1969).
Biochemical systems analysis I. Some mathematical properties of the rate law for
the component enzymatic reactions.
\textit{Journal of Theoretical Biology} \textbf{25}: 365--369.

\bibitem{voit2013}
Voit E.O. (2013).
Biochemical systems theory: a review.
\textit{ISRN Biomathematics} \textbf{2013}: 897658.

\bibitem{schuster1999}
Schuster S. et al. (1999).
Detection of elementary flux modes in biochemical networks: a promising tool
for pathway analysis and metabolic engineering.
\textit{Trends in Biotechnology} \textbf{17}: 53--60.

\bibitem{patti2012}
Patti G.J., Yanes O., Siuzdak G. (2012).
Metabolomics: the apogee of the omics trilogy.
\textit{Nature Reviews Molecular Cell Biology} \textbf{13}: 263--269.

\bibitem{chong2019}
Chong J. et al. (2019).
MetaboAnalyst 4.0: towards more transparent and integrative metabolomics analysis.
\textit{Nucleic Acids Research} \textbf{47}: W180--W187.

\bibitem{wishart2022}
Wishart D.S. et al. (2022).
HMDB 5.0: the human metabolome database for 2022.
\textit{Nucleic Acids Research} \textbf{50}: D622--D631.

\bibitem{lewis2012}
Lewis N.E. et al. (2012).
Constraining the metabolic genotype-phenotype relationship using a phylogeny of
in silico methods.
\textit{Nature Reviews Microbiology} \textbf{10}: 291--305.

\end{thebibliography}
